\makeisititle{EVALUASI AKURASI KAIDAH HARMONI FUNGSIONAL PADA MUSIK SIMBOLIK HASIL GENERASI LSTM, CNN, DAN TRANSFORMER}{\researchername}{\researchernim}{Semester Ganjil 2026/2027}{2026}
\setcounter{page}{2}

% Approval Page
\newpage
\begin{spacing}{1.15}
\begin{center}
    HALAMAN PENGESAHAN\par
    \vspace{1cm}
    \textbf{EVALUASI AKURASI KAIDAH HARMONI FUNGSIONAL PADA MUSIK SIMBOLIK HASIL GENERASI LSTM, CNN, DAN TRANSFORMER}\par
    \vspace{0.5cm}
    Diajukan Oleh:\par
    \researchername\par
    NIM: \researchernim\par
    \vspace{1.5cm}
    Proposal ini telah disetujui sebagai syarat melaksanakan Tugas Akhir/Skripsi di\par
    \programstudy, \facultyname,\par
    \institutionname\par
    \vspace{2cm}
    
    \begin{minipage}[t]{0.45\textwidth}
        \centering
        ~\\
        Mengetahui,\par
        \vspace{3cm}
        \textbf{\advisoracademic}\par
        NIP. \advisoracademicnip
    \end{minipage}
    \hfill
    \begin{minipage}[t]{0.45\textwidth}
        \centering
        Yogyakarta, 17 Juni 2026\\
        Disetujui,\par
        \vspace{3cm}
        \textbf{\advisorthesis}\par
        NIP. \advisorthesisnip
    \end{minipage}
\end{center}
\end{spacing}

% Abstract
\newpage
\begin{spacing}{1.15}
\begin{center}
    \textbf{ABSTRAK}
\end{center}

Penelitian ini mengevaluasi akurasi output model kecerdasan buatan (AI) generatif musik simbolik terhadap kaidah harmoni fungsional Barat, dengan membandingkan tiga pendekatan arsitektur yang berbeda: Recurrent (LSTM, diwakili DeepBach), Konvolusional (CNN, diwakili Coconet), dan Attention-based (Transformer, diwakili NotaGen). Dalam ranah ilmu komputer dan Music Information Retrieval (MIR), kaidah harmoni fungsional Barat diposisikan sebagai deterministic constraint, yaitu batasan biner yang dapat diverifikasi secara algoritmik, bukan sebagai standar estetika seni. Posisi ini dipakai untuk menguji sejauh mana model generatif berbasis probabilitas statistik mampu menangkap aturan yang dilarang secara absolut, bukan sekadar aturan yang jarang muncul dalam data latih. Evaluasi dilakukan secara algoritmik menggunakan pustaka music21 dengan tiga kriteria yang diformalkan dari teori harmoni fungsional Gustav Strube (1928): larangan kuint sejajar, larangan oktaf sejajar, dan resolusi nada penuntun. Hasil penelitian diharapkan menunjukkan sejauh mana pergeseran paradigma arsitektur AI, dari Recurrent dan Konvolusional generasi 2017 menuju Attention-based generasi 2025, memengaruhi constraint adherence model tersebut.\par

\textbf{Kata Kunci:} Musik Simbolik, Constraint Adherence, Harmoni Fungsional, Gustav Strube, LSTM, CNN, Transformer
\end{spacing}

% Abstract (English)
\newpage
\begin{spacing}{1.15}
\begin{center}
    \textbf{ABSTRACT}
\end{center}

This study evaluates the constraint adherence of symbolic music generative Artificial Intelligence (AI) models against Western functional harmony rules, comparing three distinct architectural paradigms: Recurrent (LSTM, represented by DeepBach), Convolutional (CNN, represented by Coconet), and Attention-based (Transformer, represented by NotaGen). Within computer science and Music Information Retrieval (MIR), Western functional harmony rules are reframed as deterministic constraints, binary in nature and algorithmically verifiable, rather than aesthetic dogma. This framing tests the extent to which probability-based generative models capture rules that are absolutely forbidden, as opposed to rules that are merely rare in training data. Evaluation is performed algorithmically using the music21 library, with three criteria formalized from Gustav Strube's (1928) theory of functional harmony: parallel fifths, parallel octaves, and leading tone resolution. Results are expected to show the extent to which the architectural paradigm shift, from Recurrent and Convolutional models in 2017 to Attention-based models in 2025, affects constraint adherence.\par

\textbf{Keywords:} Symbolic Music, Constraint Adherence, Functional Harmony, Gustav Strube, LSTM, CNN, Transformer
\end{spacing}
\newpage
